\mychapter{Referencial Teórico}{cap:tec}\lhead{REFERENCIAL TEÓRICO}

A seleção de políticas de reposição exige distinguir três objetos: o processo de demanda, a regra que produz pedidos e o critério usado para comparar seus resultados. Este capítulo define esses objetos e apresenta os métodos de previsão e inferência utilizados no trabalho. A descrição operacional dos experimentos fica no Capítulo~\ref{cap:metod}.

\section{Formulação do Problema de Reposição}
\label{sec:formulacao_problema}
\subsection{Política de Reposição: Estado, Ação e Parâmetros}
\label{sec:estado_acao_parametros}

Uma política $\pi$ transforma o estado observado $\mathbf{s}_t$ em uma quantidade de pedido não negativa:
\[
\pi:\mathcal S\rightarrow\mathcal A,\qquad Q_t=\pi(\mathbf{s}_t;\boldsymbol\theta).
\]
O estado reúne estoque disponível, pedidos em trânsito e informações de demanda disponíveis até a decisão. Os parâmetros $\boldsymbol\theta$ podem ser limiares e lotes, no caso de regras analíticas, ou pesos de uma rede neural. Meta-heurísticas ajustam parâmetros de uma regra; aprendizado por reforço ajusta a própria representação da decisão. Todas essas alternativas precisam ser avaliadas sob a mesma dinâmica de estoque.

É necessário separar estoque físico de posição de estoque. O primeiro determina o atendimento imediato; a segunda acrescenta os pedidos em trânsito e evita solicitar novamente unidades já encomendadas. Uma previsão de demanda pode integrar o estado, mas não substitui essas informações. Um vetor reduzido também não garante que o processo seja Markoviano: essa propriedade depende de representar informação suficiente para suas transições.

\subsection{Função Objetivo e Dinâmica de Estoque}
\label{sec:funcao_objetivo}

Para um horizonte finito $T$, o problema conceitual consiste em minimizar o custo esperado:
\begin{equation}
\min_\pi\mathbb E\left[\sum_{t=1}^{T}\left(hI_t+c_sS_t+c_o^{\mathrm{fix}}\mathbf1(Q_t>0)+c_o^{\mathrm{var}}Q_t\right)\right].
\label{eq:obj_inv_control}
\end{equation}
$I_t$ é o estoque após atendimento, $S_t$ a demanda perdida, $h$ o custo de manutenção por unidade e ciclo, $c_s$ o custo por unidade não atendida e $c_o^{\mathrm{fix}}$, $c_o^{\mathrm{var}}$ os custos fixo e variável de pedido. A esperança representa incerteza da demanda ou do treinamento; repetir uma trajetória determinística não estima essa distribuição.

Se $B_t$ é o estoque antes da chegada e $A_t$ a quantidade recebida no ciclo, a dinâmica de vendas perdidas é
\begin{equation}
S_t=\max(0,D_t-B_t-A_t),\quad
I_t=\max(0,B_t+A_t-D_t),\quad B_{t+1}=I_t.
\label{eq:balance}
\end{equation}
A demanda não atendida não se acumula como encomenda pendente. A convenção de chegada deve acompanhar a equação: um modelo com prazo efetivo $\ell$ usa $A_t=Q_{t-\ell}$, enquanto o simulador deste estudo usa uma fila de tamanho $L+1$, com chegada em $t+L+1$ para pedido feito em $t$. Essa diferença é explicitada na metodologia.

Minimizar custo sem restrição pode favorecer uma política que quase nunca repõe. Por isso, o trabalho acrescenta um nível mínimo de serviço. A escolha do limiar expressa um compromisso operacional e deve anteceder a comparação. Custos e serviço precisam ser apresentados conjuntamente, pois uma economia obtida com menor atendimento não caracteriza dominância.

\section{Famílias de Políticas de Reposição}
\label{sec:politicas_reposicao}
\subsection{Políticas Clássicas}
\label{sec:politicas_classicas}

As políticas clássicas oferecem referências interpretáveis. Sob demanda constante média $\mu$ por ciclo, custo fixo $K$ por pedido e manutenção $h$ por unidade e ciclo, o lote econômico é
\[
Q^*=\sqrt{2K\mu/h}.
\]
Demanda e manutenção devem usar a mesma unidade temporal. A EOQ define o tamanho do lote; um ponto de pedido separado estabelece quando solicitá-lo. Com estoque de segurança, esse ponto pode ser $s=\mu L+z\sigma\sqrt L$, sob aproximações sobre demanda e prazo. A implementação histórica será distinguida dessa fórmula porque emprega uma demanda acumulada no horizonte.

A política $(s,S)$ observa a posição de estoque $P_t$ e pede $\max(0,S-P_t)$ quando $P_t\leq s$. O limiar inferior evita pedidos a cada pequena redução de estoque; o superior determina a recomposição. A concentração de pedidos em lotes pode amplificar a variabilidade a montante, fenômeno associado ao efeito chicote~\citep{LeePadmanabhanWhang1997Bullwhip}. A intensidade desse efeito depende dos parâmetros e da demanda, não apenas do nome da política.

O modelo do Jornaleiro determina uma quantidade de período único pelo quantil de demanda associado à razão crítica entre falta e excesso. Na aproximação normal, escreve-se $Q^*=\mu+z_\alpha\sigma$. Essa expressão é ótima somente sob as hipóteses do modelo e com $\alpha$ derivado dos custos pertinentes. Usar $z=1{,}28$ como parâmetro fixo em uma regra multiperíodo constitui uma heurística inspirada no modelo, não uma garantia de otimalidade.

Regras analíticas podem receber parâmetros específicos por série e, portanto, não pressupõem necessariamente uma única parametrização para toda a rede. Sua simplicidade é uma vantagem comparativa relevante: um método mais complexo precisa demonstrar ganho além do obtido por calibração adequada de uma referência simples.

\subsection{Parametrização de Políticas via Meta-Heurísticas}

Meta-heurísticas procuram vetores $\boldsymbol\theta$ em um espaço de busca. Cada candidato induz uma política, cuja execução simulada fornece custo e serviço. Essa abordagem dispensa uma expressão diferenciável para o desempenho, mas continua sujeita a sobreajuste ao histórico usado na busca. O objetivo, os limites dos parâmetros e o orçamento de avaliações devem ser controlados na comparação.

O Algoritmo Genético (GA) mantém uma população de candidatos e alterna seleção, cruzamento e mutação~\citep{Holland1992ANAS,Talbi2009Metaheuristics}. Em reposição, cromossomos podem codificar ponto de pedido e lote; a avaliação associa cada vetor à trajetória resultante~\citep{Yuna2023IntermittentMetaheuristics}. Não há garantia geral de ótimo global em orçamento finito.

O Recozimento Simulado (SA) parte de uma solução e explora vizinhos. Uma piora $\Delta C$ pode ser aceita com probabilidade $\exp(-\Delta C/T_k)$, em que $T_k$ é a temperatura~\citep{Kirkpatrick1983SA}. Seu resfriamento controla a passagem de exploração para busca local. O resultado depende da escala da função objetivo, da vizinhança e do número de iterações.

Na Otimização por Enxame de Partículas (PSO), cada posição representa um candidato; velocidade, melhor posição individual e melhor posição coletiva orientam a atualização~\citep{Kennedy1995PSO}. A Evolução Diferencial (DE) gera candidatos pela combinação de diferenças entre vetores e aplica cruzamento e seleção~\citep{Storn1997DE}. Ambos são alternativas para espaços contínuos, mas não constituem, por si, regras operacionais distintas de reposição.

Mesmo quando duas buscas ajustam a mesma regra, seus custos computacionais e soluções podem diferir. A comparação deve identificar qual parte do ganho decorre do algoritmo de busca e qual decorre da regra parametrizada. Otimizar cada série também é diferente de selecionar, entre políticas já calibradas, qual deve ser aplicada a uma nova série.

\subsection{Políticas de Reposição via Aprendizado por Reforço}

O aprendizado por reforço representa o controle como um processo de decisão $\langle\mathcal S,\mathcal A,P,R,\gamma\rangle$. O agente escolhe pedidos, recebe uma recompensa e observa a transição do inventário. Com recompensa $r_t=-C_t$, busca maximizar o retorno descontado $\sum_t\gamma^t r_t$~\citep{SuttonBarto2018RL}. Uma restrição de serviço não é automaticamente satisfeita por essa recompensa: precisa de tratamento explícito ou verificação posterior.

DQN aproxima a função de valor $Q_\theta(s,a)$ por uma rede neural, usando memória de experiências e rede-alvo~\citep{Mnih2015DQN}. Sua aplicação requer ações discretizadas. A resolução dessa discretização limita os lotes disponíveis; episódios adicionais sobre a mesma série aumentam interações simuladas, mas não acrescentam novas realizações independentes de demanda.

PPO atualiza uma política por gradiente com uma razão de probabilidades limitada por \textit{clipping}~\citep{Schulman2017PPO}. Pode ser formulado com ações contínuas ou discretas; essa propriedade precisa ser informada para cada implementação. Comportamentos como pedidos excessivos ou ausência de reposição podem resultar da recompensa, da escala dos estados, da exploração ou do treinamento. A observação de uma trajetória degenerada não identifica isoladamente sua causa.

SARSA atualiza valores seguindo a própria política:
\[
Q(s,a)\leftarrow Q(s,a)+\eta[r+\gamma Q(s',a')-Q(s,a)].
\]
A versão tabular exige discretizar o estado; Expected SARSA substitui o próximo valor pela esperança sobre ações. A adequação de cada representação depende da informação preservada e do tamanho do problema, e não permite concluir antecipadamente que uma família será superior nas séries estudadas.

\subsection{Políticas Híbridas GA-RL}
\label{sec:ga_rl_arch}

Uma arquitetura híbrida usa a busca evolutiva para obter parâmetros ou experiências iniciais e depois ajusta uma política por reforço. No GA--DQN, transições geradas com a política do GA podem preencher a memória do agente antes do refinamento. Trata-se de uma hipótese de melhoria da inicialização; seu benefício exige comparação com os componentes isolados e controle dos respectivos orçamentos.



A inicialização não elimina o risco de sobreajuste. Se busca e aprendizado observarem a mesma trajetória usada na avaliação, o resultado descreve adaptação ao histórico, não generalização. Além disso, avaliações completas de candidatos do GA e transições armazenadas no DQN são quantidades diferentes. O trabalho trata o híbrido como candidato do portfólio; o problema de pesquisa é a seleção contextual entre candidatos.

\section{Demanda Intermitente}

Uma série intermitente alterna ciclos sem demanda e ocorrências positivas, cujos tamanhos também podem variar~\citep{Croston1972ForecastingIntermittentDemand,Boylan2008Intermittent}. Frequência e magnitude são dimensões distintas. Para $D_t\geq0$, definem-se
\begin{align}
n_+&=|\{t:D_t>0\}|,\label{eq:n_pos}\\
\mu&=T^{-1}\sum_tD_t,\qquad \sigma^2=(T-1)^{-1}\sum_t(D_t-\mu)^2.\label{eq:mu_sigma}
\end{align}
A média sobre todos os ciclos descreve o volume por período, mas pode ocultar longos intervalos sem venda. Uma sequência com valores positivos estáveis e muitos zeros é operacionalmente diferente de outra igualmente esparsa, porém com grandes picos.

\subsection[Taxonomia ADI e CV quadrado]{Taxonomia ADI e CV$^2$}
\label{sec:sb_taxa}

A classificação de \citet{Syntetos2005Categorization} combina intervalo médio entre demandas e variação dos tamanhos positivos:
\begin{align}
\mathrm{ADI}&=T/n_+,\label{eq:adi}\\
\mathrm{CV}_+^2&=(\sigma_+/\mu_+)^2.\label{eq:cv2}
\end{align}
Aqui, $\mu_+$ e $\sigma_+$ são calculados apenas nas observações positivas, diferentemente da média e variância da Equação~\ref{eq:mu_sigma}. Os limiares 1,32 e 0,49 delimitam quatro categorias, resumidas na Tabela~\ref{tab:sb_quad}~\citep{Syntetos2009InventoryForecastingReview}.

\begin{table}[htbp]
\centering
\small
\renewcommand{\arraystretch}{1.2}
\caption{Quadrantes de Syntetos-Boylan: limiares e interpretações.}
\label{tab:sb_quad}
\begin{tabular}{llll}
\toprule
Regime & ADI & CV$^2$ & Interpretação operacional \\
\midrule
\textit{Smooth} (Regular) & $< 1{,}32$ & $< 0{,}49$ & Alta frequência, demanda estável. \\
\textit{Erratic} (Errático) & $< 1{,}32$ & $\geq 0{,}49$ & Alta frequência, magnitude irregular. \\
\textit{Intermittent} (Intermitente) & $\geq 1{,}32$ & $< 0{,}49$ & Muitos zeros, magnitude relativamente estável. \\
\textit{Lumpy} (Grumoso) & $\geq 1{,}32$ & $\geq 0{,}49$ & Muitos zeros e magnitude altamente irregular. \\
\bottomrule
\end{tabular}
\end{table}

A implementação disponível também calcula uma medida de CV$^2$ incluindo zeros. Essa medida é útil como característica, mas não reproduz exatamente a definição original dos tamanhos positivos. As categorias dos artefatos são mantidas como rótulos operacionais, com essa ressalva; uma comparação entre taxonomias exige recalcular explicitamente ambos os índices.

\subsection{Efeitos na Reposição}

Demanda esparsa dificulta estimar taxas de ocorrência; tamanhos irregulares ampliam o custo de errar lotes e estoques de segurança. Croston separa a estimação dos tamanhos positivos e dos intervalos entre ocorrências. Isso pode melhorar a previsão, mas não estabelece qual regra de reposição minimiza custo sob um prazo e uma restrição de serviço específicos.



Os quadrantes fornecem uma descrição inicial, não uma prescrição de política. Séries no mesmo quadrante podem diferir em volume, tendências, sequências de zeros e concentração de picos. A hipótese do AIPE é que essas diferenças adicionais possam auxiliar a escolha; sua utilidade precisa ser medida contra uma política global bem calibrada, não inferida da classificação.

\section{Métricas de Previsão}

Erros absolutos e quadráticos dependem da escala da série; o MAPE é indefinido quando a demanda observada é zero. O estudo emprega métricas escaladas e sMAPE, com convenções explícitas para denominadores nulos~\citep{Hyndman2006MASE,Makridakis2020M4Competition}. Nenhuma métrica de previsão substitui a avaliação do efeito da decisão sobre custo e serviço.

\subsection{Métricas Utilizadas}

Sejam $e_t=D_t-\hat D_t$ os erros em $m$ pontos de avaliação e $d_1,\ldots,d_n$ o treinamento. O MASE compara o erro absoluto com a variação da referência ingênua:
\begin{equation}
\mathrm{MASE}=\frac{m^{-1}\sum_t|e_t|}{(n-1)^{-1}\sum_{j=2}^{n}|d_j-d_{j-1}|}.
\label{eq:mase}
\end{equation}
Valores menores que um indicam erro inferior à escala ingênua~\citep{Hyndman2006MASE}. A escala deve vir do treino; séries constantes requerem registro de indefinição ou regularização declarada. A presença de zeros, isoladamente, não impede o cálculo.

O sMAPE usa demanda e previsão simetricamente~\citep{Makridakis1993SMAPE}:
\begin{equation}
\mathrm{sMAPE}=\frac{200}{m}\sum_t\frac{|D_t-\hat D_t|}{|D_t|+|\hat D_t|}\,\%.
\label{eq:smape}
\end{equation}
Adota-se contribuição zero quando ambos os termos do denominador são nulos. Seu intervalo é de zero a 200\%, mas pequenas demandas podem produzir valores elevados; por isso é interpretado junto ao MASE.

A versão quadrática escalada é
\begin{equation}
\mathrm{RMSSE}=\sqrt{\frac{m^{-1}\sum_te_t^2}{(n-1)^{-1}\sum_{j=2}^{n}(d_j-d_{j-1})^2}}.
\label{eq:rmsse}
\end{equation}
Ela dá maior peso a erros grandes e também requer escala de treino não nula. O RMSSE complementa o diagnóstico; o MASE permanece a referência principal de comparação de previsão no estudo.

O viés médio, $\mathrm{MBE}=\overline D-\overline{\hat D}$, diferencia subestimação de superestimação. O U de Theil compara RMSE do modelo e da previsão ingênua~\citep{Theil1966AppliedEconometrics}. Esses indicadores auxiliam a interpretação, sem determinar a política vencedora. Uma relação entre menor erro e melhor reposição é hipótese empírica: depende da assimetria de custos, do prazo e da regra de decisão.

\section{Testes Estatísticos}
\label{sec:stat_tests}

A comparação usa observações pareadas porque as políticas atendem às mesmas séries. Testes não paramétricos evitam exigir normalidade dos KPIs, mas não dispensam atenção à independência das unidades. Produtos da mesma loja e lojas expostas às mesmas condições podem ser dependentes. Repetições idênticas não aumentam o tamanho efetivo da amostra.

\subsection{Teste de Wilcoxon}

Para diferenças $\Delta_i=X_i^A-X_i^B$, a soma de postos com sinal é
\begin{equation}
W=\sum_{i:\Delta_i\neq0}\operatorname{sgn}(\Delta_i)R_i,
\label{eq:wilcoxon}
\end{equation}
em que $R_i$ é o posto de $|\Delta_i|$. O teste de \citet{Wilcoxon1945SignedRank} avalia uma hipótese de localização sob condições sobre a distribuição das diferenças. Sua leitura como teste da mediana requer, em particular, simetria. Empates e diferenças nulas devem receber tratamento documentado.

\subsection{Teste de Friedman e Nemenyi}

Com $K$ políticas avaliadas nas mesmas $N$ instâncias, o teste de \citet{Friedman1940Ranks} compara postos médios:
\begin{equation}
\chi_F^2=\frac{12N}{K(K+1)}\sum_{j=1}^{K}\left(\bar r_j-\frac{K+1}{2}\right)^2.
\label{eq:friedman}
\end{equation}
A rejeição não identifica sozinha os pares diferentes. O procedimento de Nemenyi usa uma diferença crítica~\citep{Nemenyi1963PostHoc}:
\begin{equation}
\mathrm{DC}=q_\alpha\sqrt{\frac{K(K+1)}{6N}}.
\label{eq:cd}
\end{equation}
Postos separados por mais que DC são distinguidos sob as condições do procedimento. Um posto favorável em custo não representa necessariamente serviço aceitável; as métricas e os limiares precisam permanecer visíveis.

\subsection{Tamanho de Efeito}

A diferença padronizada de médias é reportada como descrição complementar~\citep{Cohen1988StatisticalPower}:
\begin{equation}
d=\frac{\bar X_A-\bar X_B}{s_{\mathrm{pooled}}},\qquad
s_{\mathrm{pooled}}=\sqrt{\frac{(N_A-1)s_A^2+(N_B-1)s_B^2}{N_A+N_B-2}}.
\label{eq:cohens_d}
\end{equation}
Essa forma não utiliza a covariância do pareamento; medidas baseadas nas diferenças podem ser mais apropriadas para inferência pareada. Faixas convencionais de efeito pequeno, médio ou grande não substituem a diferença de custo em unidades operacionais e sua repercussão no serviço.

O conjunto de testes, efeitos e gráficos descreve a heterogeneidade observada. Não transforma uma escolha retrospectiva de vencedor em evidência de generalização. Essa distinção orienta o protocolo de validação do próximo capítulo: a avaliação de uma política e a avaliação do mecanismo que a seleciona são etapas diferentes.
